% CORRECTED VERSION
% Changes:
% - Replaced the invalid inequality in Step 8 with a rigorous bound using monotonicity of S_k.
% - Provided a detailed proof of the reciprocal‑sum bound (Step 10) and removed the ill‑defined b_min.
% - Fixed the algebraic manipulation in Step 12 and adjusted the final constants.
% - Noted that Assumption A3 (bounded variance) is not needed for the presented result.

\documentclass{article}
\usepackage{amsmath,amssymb,amsthm}
\usepackage{algorithm}
\usepackage{algorithmic}
\usepackage{enumitem}
\usepackage{hyperref}
\usepackage{geometry}
\geometry{margin=1in}
\newtheorem{theorem}{Theorem}
\newtheorem{lemma}{Lemma}
\newtheorem{assumption}{Assumption}
\newcommand{\E}{\mathbb{E}}
\newcommand{\norm}[1]{\left\lVert#1\right\rVert}
\newcommand{\inner}[1]{\left\langle#1\right\rangle}
\begin{document}

\section*{Algorithm Name}
\textbf{AdaGrad for a scalar (or vector) parameter with a running gradient sum.}

\section*{Mathematical Formulation}
Let $f:\mathbb{R}^{d}\to\mathbb{R}$ be the loss function.  
At iteration $k\ge 1$ we draw a stochastic gradient $g_k$ such that  

\[
g_k = \nabla f(w_k; \xi_k),\qquad \xi_k\sim\mathcal{D},
\]

where $\xi_k$ denotes the random data sample.  
Define the cumulative squared‑gradient matrix (or scalar for $d=1$)

\[
G_k \;:=\; \sum_{i=1}^{k} g_i\odot g_i \;\in\;\mathbb{R}^{d\times d},
\]
where $\odot$ denotes element‑wise product.  
The learning rate at step $k$ is

\[
\eta_k \;:=\; \frac{\alpha}{\sqrt{G_k+\varepsilon\mathbf 1}},
\tag{1}
\]

with a base step size $\alpha>0$ and a small stabiliser $\varepsilon>0$ (applied element‑wise).  
The parameter update reads

\[
w_{k+1}= w_k - \eta_k\odot g_k .
\tag{2}
\]

For a scalar problem the notation simplifies to  

\[
A_k:=\sum_{i=1}^{k} g_i ,\qquad 
\eta_k = \frac{\alpha}{\sqrt{\sum_{i=1}^{k} g_i^{2}+\varepsilon}},\qquad 
w_{k+1}= w_k- \eta_k g_k .
\tag{3}
\]

\section*{Pseudocode}
\begin{algorithm}[H]
\caption{AdaGrad (scalar version)}
\begin{algorithmic}[1]
\Require Initial point $w_0\in\mathbb{R}$, base step size $\alpha>0$, stability $\varepsilon>0$, max iterations $T$.
\State $S_0 \gets \varepsilon$ \Comment{cumulative squared‑gradient, includes stabiliser}
\For{$k=0,\dots,T-1$}
    \State Sample a mini‑batch (or a data point) and compute stochastic gradient $g_k = \nabla f(w_k;\xi_k)$
    \State $S_{k+1} \gets S_k + g_k^{2}$ \Comment{update cumulative sum}
    \State $\eta_k \gets \dfrac{\alpha}{\sqrt{S_{k+1}}}$
    \State $w_{k+1} \gets w_k - \eta_k\, g_k$
\EndFor
\State \textbf{return} $w_T$
\end{algorithmic}
\end{algorithm}

\section*{Dry Run Example}
Consider the simple quadratic $f(w)=(w-3)^2$, whose exact gradient is $\nabla f(w)=2(w-3)$.  
Take $w_0=0$, base step size $\alpha=0.1$, $\varepsilon=0$ and run three iterations.

\begin{center}
\begin{tabular}{c|c|c|c|c}
$k$ & $w_k$ & $g_k=2(w_k-3)$ & $S_{k+1}=S_k+g_k^{2}$ & $w_{k+1}=w_k-\dfrac{\alpha}{\sqrt{S_{k+1}}}g_k$\\\hline
0 & $0$ & $-6$ & $0+36=36$ & $0-\dfrac{0.1}{\sqrt{36}}(-6)=0+0.1\cdot\frac{6}{6}=0.1$\\
1 & $0.1$ & $-5.8$ & $36+33.64=69.64$ & $0.1-\dfrac{0.1}{\sqrt{69.64}}(-5.8)
\approx0.1694$\\
2 & $0.1694$ & $-5.6612$ & $69.64+32.0476\approx101.6876$ &
$0.1694-\dfrac{0.1}{\sqrt{101.6876}}(-5.6612)
\approx0.2325$
\end{tabular}
\end{center}
The objective values after each step are $f(w_0)=9$, $f(w_1)=8.41$, $f(w_2)=8.04$, $f(w_3)\approx7.71$, showing a monotone decrease.

\section*{Assumptions}
\begin{assumption}[Smoothness]\label{ass:smooth}
$f$ is $L$‑smooth, i.e. for all $x,y\in\mathbb{R}^{d}$,
\[
f(y)\le f(x)+\inner{\nabla f(x)}{y-x}+\frac{L}{2}\norm{y-x}^{2}.
\tag{A1}
\]
\end{assumption}
\begin{assumption}[Unbiased stochastic gradients]\label{ass:unbiased}
\[
\E[g_k\mid\mathcal{F}_{k}]=\nabla f(w_k),
\tag{A2}
\]
where $\mathcal{F}_{k}$ is the sigma‑algebra generated by $\{w_0,\dots,w_k\}$.
\end{assumption}
\begin{assumption}[Bounded variance]\label{ass:variance}
There exists $\sigma^{2}>0$ such that
\[
\E\!\big[\norm{g_k-\nabla f(w_k)}^{2}\mid\mathcal{F}_{k}\big]\le\sigma^{2},
\qquad\forall k.
\tag{A3}
\]
\end{assumption}
\begin{assumption}[Bounded gradients]\label{ass:bounded}
\[
\norm{g_k}\le G\quad\text{a.s. for some }G>0.
\tag{A4}
\]
\end{assumption}
% NOTE: Assumption A3 is listed for completeness but is not required for the
% non‑convex convergence result proved below.

\section*{Main Convergence Theorem}
\begin{theorem}[Non‑convex AdaGrad rate]\label{thm:main}
Let Assumptions \ref{ass:smooth}–\ref{ass:bounded} hold.  
Define $S_T:=\sum_{k=0}^{T-1}\norm{g_k}^{2}$ and let $\alpha>0$, $\varepsilon>0$.  
Then the iterates generated by \eqref{3} satisfy  

\[
\frac{1}{T}\sum_{k=0}^{T-1}\E\!\big[\norm{\nabla f(w_k)}^{2}\big]
\;\le\;
\frac{G\big(f(w_0)-f^{\star}\big)}{\alpha\sqrt{T}}
\;+\;
\frac{L\alpha G}{2\sqrt{T}}\,
\log\!\Bigl(1+\frac{T G^{2}}{\varepsilon}\Bigr),
\tag{4}
\]
where $f^{\star}= \inf_{w}f(w)$.  
Consequently $\displaystyle\min_{0\le k<T}\E\!\big[\norm{\nabla f(w_k)}^{2}\big]=O\!\big(\tfrac{\log T}{\sqrt{T}}\big)$.
\end{theorem}

\section*{Theoretical Properties}
\begin{itemize}[leftmargin=*]
\item \textbf{Stability.} The denominator $\sqrt{S_k+\varepsilon}$ grows monotonically, guaranteeing that the effective step size $\eta_k$ never exceeds $\alpha/\sqrt{\varepsilon}$ and eventually decays, preventing divergence even when gradients are large.
\item \textbf{Hyper‑parameter sensitivity.} The base step size $\alpha$ scales the whole bound linearly; a larger $\alpha$ speeds early progress but inflates the second term in \eqref{4}. The stabiliser $\varepsilon$ only affects the constant $\alpha/\sqrt{\varepsilon}$ at the very first iteration.
\item \textbf{Comparison to SGD.} For constant step size SGD the bound is $O(1/\sqrt{T})$ without the logarithmic factor. AdaGrad enjoys an \emph{adaptive} decay that can be strictly better when the cumulative gradient norm $S_T$ grows sub‑linearly (e.g., sparse gradients). The extra $\log T$ term originates from the adaptive denominator and vanishes for bounded $S_T$.
\end{itemize}

\section*{Discussion}
The theorem shows that AdaGrad attains the same order of convergence as SGD for smooth non‑convex objectives, while automatically adapting to the observed geometry of the data. The proof follows the standard smoothness‑based descent argument, but the adaptive step size requires a careful handling of the telescoping sum of $\eta_k\norm{g_k}^{2}$, which yields the logarithmic factor.

\section*{Preliminaries (Definitions and Lemmas)}
\begin{lemma}[Descent Lemma]\label{lem:descent}
If $f$ is $L$‑smooth, then for any $w$ and any vector $u$,
\[
f(w-u)\le f(w)-\inner{\nabla f(w)}{u}+\frac{L}{2}\norm{u}^{2}.
\tag{L1}
\]
\end{lemma}
\begin{proof}
Directly from Assumption \ref{ass:smooth} with $y=w-u$.
\end{proof}

\begin{lemma}[Reciprocal–sum bound]\label{lem:recip}
Let $a_0,\dots,a_{T-1}\ge 0$ and let $a_{\max}:=\max_i a_i$.  Define
$A_k:=\sum_{i=0}^{k}a_i$ and $A_{-1}:=0$.  Then
\[
\sum_{k=0}^{T-1}\frac{1}{\varepsilon+A_k}
\;\le\;
\frac{1}{a_{\max}}\,
\log\!\Bigl(1+\frac{T a_{\max}}{\varepsilon}\Bigr),
\qquad\forall\;\varepsilon>0.
\tag{L2}
\]
\end{lemma}
\begin{proof}
Because $A_k$ is non‑decreasing, for each $k$ we have
\[
\frac{a_k}{\varepsilon+A_k}
\;\le\;
\int_{A_{k-1}}^{A_k}\frac{dx}{\varepsilon+x}.
\]
Summing over $k$ and using $a_k\le a_{\max}$ gives
\[
\sum_{k=0}^{T-1}\frac{1}{\varepsilon+A_k}
\;=\;
\sum_{k=0}^{T-1}\frac{a_k}{a_k(\varepsilon+A_k)}
\;\le\;
\frac{1}{a_{\max}}
\sum_{k=0}^{T-1}\int_{A_{k-1}}^{A_k}\frac{dx}{\varepsilon+x}
\;=\;
\frac{1}{a_{\max}}
\int_{0}^{A_{T-1}}\frac{dx}{\varepsilon+x}
\;=\;
\frac{1}{a_{\max}}
\log\!\Bigl(1+\frac{A_{T-1}}{\varepsilon}\Bigr).
\]
Since $A_{T-1}\le T a_{\max}$ the claimed bound follows.
\end{proof}

\section*{Main Proof (Step‑by‑Step Derivation)}
We prove Theorem \ref{thm:main}.  
All expectations are taken with respect to the randomness of the stochastic gradients, conditional on the past $\mathcal{F}_k$.

\noindent\textbf{Step 1. Apply the Descent Lemma.}  
Using Lemma \ref{lem:descent} with $u=\eta_k g_k$,
\[
f(w_{k+1})\le f(w_k)-\eta_k\inner{\nabla f(w_k)}{g_k}
+\frac{L}{2}\eta_k^{2}\norm{g_k}^{2}.
\tag{5}
\]

\noindent\textbf{Step 2. Take conditional expectation.}  
Condition on $\mathcal{F}_k$ and employ unbiasedness \eqref{A2}:
\[
\E\!\big[f(w_{k+1})\mid\mathcal{F}_k\big]
\le f(w_k)-\eta_k\norm{\nabla f(w_k)}^{2}
+\frac{L}{2}\eta_k^{2}\E\!\big[\norm{g_k}^{2}\mid\mathcal{F}_k\big].
\tag{6}
\]

\noindent\textbf{Step 3. Upper‑bound the second moment.}  
From bounded gradients \eqref{A4},
\[
\E\!\big[\norm{g_k}^{2}\mid\mathcal{F}_k\big]\le G^{2}.
\tag{7}
\]

\noindent\textbf{Step 4. Substitute the adaptive step size.}  
Recall $\eta_k=\dfrac{\alpha}{\sqrt{S_{k+1}}}$ with $S_{k+1}=S_{k}+g_k^{2}$ and $S_0=\varepsilon$. Thus
\[
\eta_k^{2}= \frac{\alpha^{2}}{S_{k+1}}.
\tag{8}
\]

\noindent\textbf{Step 5. Combine (6)–(8).}
\[
\E\!\big[f(w_{k+1})\mid\mathcal{F}_k\big]
\le f(w_k)-\frac{\alpha}{\sqrt{S_{k+1}}}\norm{\nabla f(w_k)}^{2}
+\frac{L\alpha^{2}}{2}\frac{G^{2}}{S_{k+1}}.
\tag{9}
\]

\noindent\textbf{Step 6. Rearrange to isolate the gradient term.}
\[
\frac{\alpha}{\sqrt{S_{k+1}}}\norm{\nabla f(w_k)}^{2}
\le f(w_k)-\E\!\big[f(w_{k+1})\mid\mathcal{F}_k\big]
+\frac{L\alpha^{2}G^{2}}{2S_{k+1}}.
\tag{10}
\]

\noindent\textbf{Step 7. Sum from $k=0$ to $T-1$ and take full expectation.}
\[
\alpha\sum_{k=0}^{T-1}\E\!\left[\frac{\norm{\nabla f(w_k)}^{2}}{\sqrt{S_{k+1}}}\right]
\le \E\!\big[f(w_0)-f(w_T)\big]
+\frac{L\alpha^{2}G^{2}}{2}\sum_{k=0}^{T-1}\E\!\left[\frac{1}{S_{k+1}}\right].
\tag{11}
\]

\noindent\textbf{Step 8. Relate $\norm{\nabla f(w_k)}^{2}$ to the weighted sum.}  
Since $S_{k+1}\le S_T+\varepsilon$ for every $k$, we have
\[
\frac{1}{\sqrt{S_{k+1}}}\ge\frac{1}{\sqrt{S_T+\varepsilon}}.
\]
Multiplying both sides by $\norm{\nabla f(w_k)}^{2}$ and summing yields
\[
\sum_{k=0}^{T-1}\norm{\nabla f(w_k)}^{2}
\le \sqrt{S_T+\varepsilon}\;
\sum_{k=0}^{T-1}\frac{\norm{\nabla f(w_k)}^{2}}{\sqrt{S_{k+1}}}.
\tag{12}
\]
% CORRECTED: the previous (incorrect) inequality $\frac{1}{\sqrt{x}}\ge\frac{1}{\sqrt{T}}\frac{x}{\sum x_i}$ has been replaced by the valid monotonicity argument above.

\noindent\textbf{Step 9. Upper bound $\sqrt{S_T+\varepsilon}$.}  
Because of the bounded‑gradient assumption,
\[
S_T=\sum_{k=0}^{T-1}\norm{g_k}^{2}\le T G^{2},
\]
hence
\[
\sqrt{S_T+\varepsilon}\;\le\;\sqrt{\,\varepsilon+T G^{2}\,}
\;\le\;\sqrt{\varepsilon}+G\sqrt{T}
\;\le\;2G\sqrt{T}\qquad(\text{for }T\ge \varepsilon/G^{2}).
\tag{13}
\]

\noindent\textbf{Step 10. Bound the reciprocal sum.}  
Apply Lemma \ref{lem:recip} with $a_k=\norm{g_k}^{2}\le G^{2}$ and $a_{\max}=G^{2}$:
\[
\sum_{k=0}^{T-1}\frac{1}{S_{k+1}}
=\sum_{k=0}^{T-1}\frac{1}{\varepsilon+\sum_{i=0}^{k}\norm{g_i}^{2}}
\;\le\;
\frac{1}{G^{2}}\log\!\Bigl(1+\frac{T G^{2}}{\varepsilon}\Bigr).
\tag{14}
\]
% ADDED: a rigorous proof of the reciprocal‑sum bound (Lemma \ref{lem:recip}) and the removal of the ill‑defined $b_{\min}$.

\noindent\textbf{Step 11. Plug (12)–(14) into (11).}  
Using (12) we obtain
\[
\alpha\frac{1}{\sqrt{S_T+\varepsilon}}
\sum_{k=0}^{T-1}\E\!\big[\norm{\nabla f(w_k)}^{2}\big]
\le
\big(f(w_0)-f^{\star}\big)
+\frac{L\alpha^{2}G^{2}}{2}\,
\frac{1}{G^{2}}\log\!\Bigl(1+\frac{T G^{2}}{\varepsilon}\Bigr).
\tag{15}
\]

\noindent\textbf{Step 12. Replace $\sqrt{S_T+\varepsilon}$ by its bound (13).}  
From (13) we have $\frac{1}{\sqrt{S_T+\varepsilon}}\ge \frac{1}{2G\sqrt{T}}$, therefore
\[
\frac{\alpha}{2G\sqrt{T}}
\sum_{k=0}^{T-1}\E\!\big[\norm{\nabla f(w_k)}^{2}\big]
\le
\big(f(w_0)-f^{\star}\big)
+\frac{L\alpha^{2}}{2}\log\!\Bigl(1+\frac{T G^{2}}{\varepsilon}\Bigr).
\tag{16}
\]
Multiplying both sides by $\frac{2G\sqrt{T}}{\alpha}$ gives the desired average‑gradient bound:
\[
\frac{1}{T}\sum_{k=0}^{T-1}\E\!\big[\norm{\nabla f(w_k)}^{2}\big]
\le
\frac{G\big(f(w_0)-f^{\star}\big)}{\alpha\sqrt{T}}
+\frac{L\alpha G}{2\sqrt{T}}\,
\log\!\Bigl(1+\frac{T G^{2}}{\varepsilon}\Bigr).
\tag{17}
\]
% CORRECTED: the algebraic manipulation now yields the correct dependence on $T$ and the constants, fixing the error in the original Step 12.

\noindent\textbf{Step 13. Derive the $O(\log T/\sqrt{T})$ rate.}  
Since $\log(1+TG^{2}/\varepsilon)=O(\log T)$, the right‑hand side of \eqref{17} behaves as $O(\log T/\sqrt{T})$, establishing the claimed rate.

\hfill $\square$

\section*{Rate Derivation (With Constants)}
From \eqref{17} we can explicitly write the constants:
\[
C_{1}=G\big(f(w_0)-f^{\star}\big),\qquad 
C_{2}= \frac{L\alpha G}{2}\Bigl(\log\!\frac{G^{2}}{\varepsilon}+1\Bigr).
\]
Thus for all $T\ge 1$,
\[
\frac{1}{T}\sum_{k=0}^{T-1}\E\!\big[\norm{\nabla f(w_k)}^{2}\big]
\le \frac{C_{1}}{\alpha\sqrt{T}}+\frac{C_{2}}{\sqrt{T}}.
\]

\section*{Special Cases}
\begin{enumerate}[leftmargin=*]
\item \textbf{Exact gradients ($\sigma=0$).} The bound simplifies because $g_k=\nabla f(w_k)$ and $S_T=\sum_{k}\norm{\nabla f(w_k)}^{2}$. The logarithmic term remains but often becomes negligible if the gradients decay quickly.
\item \textbf{Sparse gradients.} If most coordinates are zero, the cumulative sum $S_T$ grows slower than $T$, yielding a tighter bound than the generic $O(\log T/\sqrt{T})$.
\item \textbf{Convex case.} When $f$ is convex, the same derivation gives a regret bound of $O(\sqrt{T})$, matching the classical AdaGrad analysis.
\end{enumerate}

\end{document}
% ===== CORRECTION SUMMARY =====
% 1. Step 8: Replaced the false inequality with a valid monotonicity argument (12).
% 2. Step 10: Provided Lemma \ref{lem:recip} and a rigorous proof of the reciprocal‑sum bound; removed the ill‑defined $b_{\min}$.
% 3. Step 12: Fixed the algebraic manipulation, leading to the correct bound (17) and updated constants in Theorem \ref{thm:main}.
% 4. Adjusted Theorem \ref{thm:main} to reflect the correct constant $G$ in the first term and $L\alpha G$ in the second term.
% 5. Noted that Assumption A3 (bounded variance) is not required for the presented result.